\documentclass[aspectratio=169,11pt]{beamer}
\usetheme{Madrid}
\usecolortheme{default}
\usepackage[T1]{fontenc}
\setbeamertemplate{navigation symbols}{}
\setbeamertemplate{footline}[frame number]
\usepackage{enumitem}
\usepackage{amsmath,amssymb}
\usepackage{booktabs}
\usepackage{array}

\title{Databricks ML Associate}
\subtitle{High-Yield Next 100 (Set 2)}
\author{Practice Exam Deck}
\date{}

\begin{document}

\begin{frame}
\titlepage
\end{frame}

\begin{frame}{How to Use This Deck}
\begin{itemize}
\item 100 more high-yield questions (Q1--Q100), complementing Set 1, on topics that commonly appear on the Databricks Machine Learning Associate exam.
\item Each question slide has four options (A--D); the next slide gives the correct answer and why the others are wrong.
\item Topic order: platform, compute, Delta and Spark; MLflow; model registry; Feature Engineering; data processing; tuning; Spark ML; metrics and concepts; deployment, jobs, and monitoring.
\end{itemize}
\end{frame}

\begin{frame}{Exam Domains Covered}
\begin{itemize}
\item \textbf{Databricks Machine Learning:} runtime, AutoML, MLflow, Unity Catalog registry, Feature Engineering client.
\item \textbf{Data Processing:} EDA, imputation, encoding, Spark ML transformers, pandas API on Spark and UDFs.
\item \textbf{Model Development:} Hyperopt, CrossValidator, metrics, bias/variance, leakage.
\item \textbf{Model Deployment:} batch, streaming, real-time serving, monitoring.
\end{itemize}
\end{frame}

% ---------------------------------------------------------------
\begin{frame}{Question 1}
What does the auto-termination setting on a cluster do?
\begin{enumerate}[label=\Alph*.]
\item Restarts the cluster nightly
\item Deletes the notebooks attached
\item Shuts the cluster down after a set period of inactivity
\item Stops a job when it fails
\end{enumerate}
\end{frame}
\begin{frame}{Answer 1}
\textbf{Correct: C} --- it saves cost on idle compute.\\
\vspace{0.5em}
\textbf{Wrong:}
\begin{itemize}
\item A: not what it does.
\item B: notebooks are not deleted.
\item D: that is a job setting.
\end{itemize}
\end{frame}

% ---------------------------------------------------------------
\begin{frame}{Question 2}
What is the purpose of a cluster policy?
\begin{enumerate}[label=\Alph*.]
\item To schedule notebooks
\item To register models
\item To store secrets
\item To restrict and standardize the configuration users can choose
\end{enumerate}
\end{frame}
\begin{frame}{Answer 2}
\textbf{Correct: D} --- admins use policies to control size, runtime, and cost.\\
\vspace{0.5em}
\textbf{Wrong:}
\begin{itemize}
\item A: that is Workflows.
\item B: that is the model registry.
\item C: that is secret scopes.
\end{itemize}
\end{frame}

% ---------------------------------------------------------------
\begin{frame}{Question 3}
What does cluster autoscaling do?
\begin{enumerate}[label=\Alph*.]
\item Changes the driver runtime version
\item Adds or removes workers between a minimum and maximum based on load
\item Increases the dataset size
\item Scales the model endpoint to zero
\end{enumerate}
\end{frame}
\begin{frame}{Answer 3}
\textbf{Correct: B} --- it matches compute to the workload.\\
\vspace{0.5em}
\textbf{Wrong:}
\begin{itemize}
\item A: not what it does.
\item C: not what it does.
\item D: that is serving.
\end{itemize}
\end{frame}

% ---------------------------------------------------------------
\begin{frame}{Question 4}
Which magic command runs another notebook inline so its functions and variables are available?
\begin{enumerate}[label=\Alph*.]
\item \texttt{\%fs}
\item \texttt{\%md}
\item \texttt{\%run}
\item \texttt{\%sh}
\end{enumerate}
\end{frame}
\begin{frame}{Answer 4}
\textbf{Correct: C} --- \texttt{\%run} includes the other notebook in the current one.\\
\vspace{0.5em}
\textbf{Wrong:}
\begin{itemize}
\item A: runs file system commands.
\item B: renders markdown.
\item D: runs shell commands.
\end{itemize}
\end{frame}

% ---------------------------------------------------------------
\begin{frame}{Question 5}
After \texttt{\%pip install} changes a package already imported, what should you do?
\begin{enumerate}[label=\Alph*.]
\item Run \texttt{VACUUM}
\item Run \texttt{dbutils.library.restartPython()}
\item Run \texttt{df.cache()}
\item Run \texttt{mlflow.end\_run()}
\end{enumerate}
\end{frame}
\begin{frame}{Answer 5}
\textbf{Correct: B} --- restarting Python picks up the new package version.\\
\vspace{0.5em}
\textbf{Wrong:}
\begin{itemize}
\item A: a Delta command.
\item C: unrelated.
\item D: ends a run only.
\end{itemize}
\end{frame}

% ---------------------------------------------------------------
\begin{frame}{Question 6}
Which Delta Lake command performs an upsert?
\begin{enumerate}[label=\Alph*.]
\item \texttt{OPTIMIZE}
\item \texttt{DESCRIBE HISTORY}
\item \texttt{MERGE INTO}
\item \texttt{INSERT OVERWRITE DIRECTORY}
\end{enumerate}
\end{frame}
\begin{frame}{Answer 6}
\textbf{Correct: C} --- it updates matching rows and inserts new ones.\\
\vspace{0.5em}
\textbf{Wrong:}
\begin{itemize}
\item A: compacts files.
\item B: shows the table log.
\item D: writes files to a directory.
\end{itemize}
\end{frame}

% ---------------------------------------------------------------
\begin{frame}{Question 7}
What does \texttt{OPTIMIZE} do on a Delta table?
\begin{enumerate}[label=\Alph*.]
\item Deletes all old versions
\item Encrypts the table
\item Rolls back to an earlier version
\item Compacts many small files into larger ones
\end{enumerate}
\end{frame}
\begin{frame}{Answer 7}
\textbf{Correct: D} --- it improves read performance.\\
\vspace{0.5em}
\textbf{Wrong:}
\begin{itemize}
\item A: that is \texttt{VACUUM}.
\item B: not what it does.
\item C: that is \texttt{RESTORE}.
\end{itemize}
\end{frame}

% ---------------------------------------------------------------
\begin{frame}{Question 8}
What does \texttt{VACUUM} do on a Delta table?
\begin{enumerate}[label=\Alph*.]
\item Clones the table
\item Removes data files no longer referenced and older than the retention period
\item Adds a column
\item Compacts small files
\end{enumerate}
\end{frame}
\begin{frame}{Answer 8}
\textbf{Correct: B} --- this frees storage but limits time travel.\\
\vspace{0.5em}
\textbf{Wrong:}
\begin{itemize}
\item A: that is \texttt{CLONE}.
\item C: that is \texttt{ALTER TABLE}.
\item D: that is \texttt{OPTIMIZE}.
\end{itemize}
\end{frame}

% ---------------------------------------------------------------
\begin{frame}{Question 9}
Which command shows the version history of a Delta table?
\begin{enumerate}[label=\Alph*.]
\item \texttt{DESCRIBE TABLE EXTENDED} only
\item \texttt{SHOW TABLES}
\item \texttt{DESCRIBE HISTORY table\_name}
\item \texttt{EXPLAIN table\_name}
\end{enumerate}
\end{frame}
\begin{frame}{Answer 9}
\textbf{Correct: C} --- it lists versions, timestamps, and operations.\\
\vspace{0.5em}
\textbf{Wrong:}
\begin{itemize}
\item A: shows metadata, not versions.
\item B: lists tables.
\item D: shows a query plan.
\end{itemize}
\end{frame}

% ---------------------------------------------------------------
\begin{frame}{Question 10}
Which command restores a Delta table to an earlier version?
\begin{enumerate}[label=\Alph*.]
\item \texttt{UNDO TABLE t 5}
\item \texttt{OPTIMIZE t VERSION 5}
\item \texttt{VACUUM t VERSION 5}
\item \texttt{RESTORE TABLE t TO VERSION AS OF 5}
\end{enumerate}
\end{frame}
\begin{frame}{Answer 10}
\textbf{Correct: D} --- it rolls the table back to that version.\\
\vspace{0.5em}
\textbf{Wrong:}
\begin{itemize}
\item A: not valid.
\item B: OPTIMIZE compacts files.
\item C: VACUUM deletes files.
\end{itemize}
\end{frame}

% ---------------------------------------------------------------
\begin{frame}{Question 11}
What does \texttt{ZORDER BY} do in \texttt{OPTIMIZE table ZORDER BY (col)}?
\begin{enumerate}[label=\Alph*.]
\item Encrypts the column
\item Co-locates related values of the column in the same files to speed up filters
\item Sorts the output of a query only
\item Deletes duplicate rows
\end{enumerate}
\end{frame}
\begin{frame}{Answer 11}
\textbf{Correct: B} --- it improves data skipping.\\
\vspace{0.5em}
\textbf{Wrong:}
\begin{itemize}
\item A: not what it does.
\item C: it changes file layout.
\item D: not what it does.
\end{itemize}
\end{frame}

% ---------------------------------------------------------------
\begin{frame}{Question 12}
Which write saves a DataFrame as a managed Delta table?
\begin{enumerate}[label=\Alph*.]
\item \texttt{df.cache("t")}
\item \texttt{df.toTable("t")}
\item \texttt{df.save\_delta("t")}
\item \texttt{df.write.mode("overwrite").saveAsTable("catalog.schema.t")}
\end{enumerate}
\end{frame}
\begin{frame}{Answer 12}
\textbf{Correct: D} --- saveAsTable registers the table in the catalog.\\
\vspace{0.5em}
\textbf{Wrong:}
\begin{itemize}
\item A: caching does not persist a table.
\item B: not an API.
\item C: not an API.
\end{itemize}
\end{frame}

% ---------------------------------------------------------------
\begin{frame}{Question 13}
What is the difference between \texttt{repartition(n)} and \texttt{coalesce(n)}?
\begin{enumerate}[label=\Alph*.]
\item They are identical
\item \texttt{coalesce} always increases partitions
\item \texttt{coalesce} reduces partitions without a full shuffle; \texttt{repartition} does a full shuffle
\item \texttt{repartition} only decreases partitions
\end{enumerate}
\end{frame}
\begin{frame}{Answer 13}
\textbf{Correct: C} --- coalesce is cheaper when decreasing partitions.\\
\vspace{0.5em}
\textbf{Wrong:}
\begin{itemize}
\item A: they differ in shuffling.
\item B: it only decreases.
\item D: it can increase or decrease.
\end{itemize}
\end{frame}

% ---------------------------------------------------------------
\begin{frame}{Question 14}
When is a broadcast join beneficial?
\begin{enumerate}[label=\Alph*.]
\item When both tables are huge
\item When using only streaming data
\item When one table is small enough to copy to every executor
\item When there is no join key
\end{enumerate}
\end{frame}
\begin{frame}{Answer 14}
\textbf{Correct: C} --- it avoids shuffling the large table.\\
\vspace{0.5em}
\textbf{Wrong:}
\begin{itemize}
\item A: broadcasting would not fit.
\item B: unrelated.
\item D: unrelated.
\end{itemize}
\end{frame}

% ---------------------------------------------------------------
\begin{frame}{Question 15}
What is a shuffle in Spark?
\begin{enumerate}[label=\Alph*.]
\item Randomly ordering the notebook cells
\item Copying a table to another catalog
\item Encrypting data in memory
\item Redistribution of data across partitions and workers, often costly
\end{enumerate}
\end{frame}
\begin{frame}{Answer 15}
\textbf{Correct: D} --- joins and groupBy usually cause shuffles.\\
\vspace{0.5em}
\textbf{Wrong:}
\begin{itemize}
\item A: unrelated.
\item B: unrelated.
\item C: unrelated.
\end{itemize}
\end{frame}

% ---------------------------------------------------------------
\begin{frame}{Question 16}
Which Spark call triggers computation?
\begin{enumerate}[label=\Alph*.]
\item \texttt{df.count()}
\item \texttt{df.select("a")}
\item \texttt{df.filter("a > 1")}
\item \texttt{df.withColumn("b", df.a)}
\end{enumerate}
\end{frame}
\begin{frame}{Answer 16}
\textbf{Correct: A} --- count is an action.\\
\vspace{0.5em}
\textbf{Wrong:}
\begin{itemize}
\item B: a transformation.
\item C: a transformation.
\item D: a transformation.
\end{itemize}
\end{frame}

% ---------------------------------------------------------------
\begin{frame}{Question 17}
Which privilege lets a user read a Unity Catalog table (along with USE CATALOG and USE SCHEMA)?
\begin{enumerate}[label=\Alph*.]
\item \texttt{SELECT}
\item \texttt{WRITE VOLUME}
\item \texttt{EXECUTE}
\item \texttt{CREATE MODEL}
\end{enumerate}
\end{frame}
\begin{frame}{Answer 17}
\textbf{Correct: A} --- SELECT grants read access to table data.\\
\vspace{0.5em}
\textbf{Wrong:}
\begin{itemize}
\item B: applies to volumes.
\item C: is for functions and models.
\item D: creates models in a schema.
\end{itemize}
\end{frame}

% ---------------------------------------------------------------
\begin{frame}{Question 18}
What does Unity Catalog lineage provide?
\begin{enumerate}[label=\Alph*.]
\item Faster clusters
\item Automatic tracking of how tables, columns, and models relate to upstream data
\item Automatic hyperparameter tuning
\item Free storage
\end{enumerate}
\end{frame}
\begin{frame}{Answer 18}
\textbf{Correct: B} --- it helps with impact analysis and governance.\\
\vspace{0.5em}
\textbf{Wrong:}
\begin{itemize}
\item A: not provided.
\item C: not provided.
\item D: not provided.
\end{itemize}
\end{frame}

% ---------------------------------------------------------------
\begin{frame}{Question 19}
Which call logs several hyperparameters at once?
\begin{enumerate}[label=\Alph*.]
\item \texttt{mlflow.log\_artifacts\_params()}
\item \texttt{mlflow.log\_params(\{"max\_depth": 5, "n\_estimators": 100\})}
\item \texttt{mlflow.set\_params()}
\item \texttt{mlflow.log\_param\_dict()}
\end{enumerate}
\end{frame}
\begin{frame}{Answer 19}
\textbf{Correct: B} --- log\_params takes a dictionary.\\
\vspace{0.5em}
\textbf{Wrong:}
\begin{itemize}
\item A: not an API.
\item C: not an API.
\item D: not an API.
\end{itemize}
\end{frame}

% ---------------------------------------------------------------
\begin{frame}{Question 20}
Which call logs several metrics at once?
\begin{enumerate}[label=\Alph*.]
\item \texttt{mlflow.log\_score()}
\item \texttt{mlflow.log\_metric\_dict()}
\item \texttt{mlflow.set\_metrics()}
\item \texttt{mlflow.log\_metrics(\{"rmse": 3.1, "r2": 0.8\})}
\end{enumerate}
\end{frame}
\begin{frame}{Answer 20}
\textbf{Correct: D} --- log\_metrics takes a dictionary.\\
\vspace{0.5em}
\textbf{Wrong:}
\begin{itemize}
\item A: not an API.
\item B: not an API.
\item C: not an API.
\end{itemize}
\end{frame}

% ---------------------------------------------------------------
\begin{frame}{Question 21}
What is \texttt{mlflow.set\_tag("team", "risk")} best used for?
\begin{enumerate}[label=\Alph*.]
\item Adding searchable string metadata to a run
\item Saving model weights
\item Logging a numeric metric
\item Changing the experiment
\end{enumerate}
\end{frame}
\begin{frame}{Answer 21}
\textbf{Correct: A} --- tags label and organize runs.\\
\vspace{0.5em}
\textbf{Wrong:}
\begin{itemize}
\item B: that is \texttt{log\_model}.
\item C: that is \texttt{log\_metric}.
\item D: that is \texttt{set\_experiment}.
\end{itemize}
\end{frame}

% ---------------------------------------------------------------
\begin{frame}{Question 22}
Which call returns the currently active run object?
\begin{enumerate}[label=\Alph*.]
\item \texttt{mlflow.current\_model()}
\item \texttt{mlflow.active\_run()}
\item \texttt{mlflow.run\_id()}
\item \texttt{mlflow.get\_active\_experiment()}
\end{enumerate}
\end{frame}
\begin{frame}{Answer 22}
\textbf{Correct: B} --- it returns the active run or None.\\
\vspace{0.5em}
\textbf{Wrong:}
\begin{itemize}
\item A: not an API.
\item C: not an API.
\item D: not the run.
\end{itemize}
\end{frame}

% ---------------------------------------------------------------
\begin{frame}{Question 23}
How do you read the run ID inside an active run?
\begin{enumerate}[label=\Alph*.]
\item \texttt{mlflow.get\_id()}
\item \texttt{mlflow.run\_id}
\item \texttt{mlflow.active\_run().metrics}
\item \texttt{mlflow.active\_run().info.run\_id}
\end{enumerate}
\end{frame}
\begin{frame}{Answer 23}
\textbf{Correct: D} --- run\_id is stored in the run info.\\
\vspace{0.5em}
\textbf{Wrong:}
\begin{itemize}
\item A: not an API.
\item B: not an attribute.
\item C: holds metrics.
\end{itemize}
\end{frame}

% ---------------------------------------------------------------
\begin{frame}{Question 24}
What does \texttt{mlflow.sklearn.load\_model(uri)} return?
\begin{enumerate}[label=\Alph*.]
\item A Spark DataFrame
\item The native scikit-learn model object
\item A run object
\item A PyFuncModel always
\end{enumerate}
\end{frame}
\begin{frame}{Answer 24}
\textbf{Correct: B} --- flavor-specific loaders return the native object.\\
\vspace{0.5em}
\textbf{Wrong:}
\begin{itemize}
\item A: not what it returns.
\item C: not what it returns.
\item D: that is \texttt{mlflow.pyfunc.load\_model}.
\end{itemize}
\end{frame}

% ---------------------------------------------------------------
\begin{frame}{Question 25}
What does \texttt{mlflow.pyfunc.load\_model(uri)} return?
\begin{enumerate}[label=\Alph*.]
\item A Hyperopt Trials object
\item A generic model with a \texttt{predict()} method
\item The raw training DataFrame
\item A feature table
\end{enumerate}
\end{frame}
\begin{frame}{Answer 25}
\textbf{Correct: B} --- it works across flavors.\\
\vspace{0.5em}
\textbf{Wrong:}
\begin{itemize}
\item A: not what it returns.
\item C: not what it returns.
\item D: not what it returns.
\end{itemize}
\end{frame}

% ---------------------------------------------------------------
\begin{frame}{Question 26}
Which file in a logged MLflow model folder describes its flavors and metadata?
\begin{enumerate}[label=\Alph*.]
\item \texttt{cluster.json}
\item \texttt{README.md}
\item \texttt{train.csv}
\item \texttt{MLmodel}
\end{enumerate}
\end{frame}
\begin{frame}{Answer 26}
\textbf{Correct: D} --- MLmodel lists flavors, signature, and more.\\
\vspace{0.5em}
\textbf{Wrong:}
\begin{itemize}
\item A: not generated.
\item B: not generated.
\item C: not generated.
\end{itemize}
\end{frame}

% ---------------------------------------------------------------
\begin{frame}{Question 27}
Which files record the model's Python dependencies?
\begin{enumerate}[label=\Alph*.]
\item \texttt{requirements.txt} and \texttt{conda.yaml}
\item \texttt{MLmodel} only
\item \texttt{notebook.ipynb}
\item \texttt{signature.json}
\end{enumerate}
\end{frame}
\begin{frame}{Answer 27}
\textbf{Correct: A} --- they let the environment be recreated for serving.\\
\vspace{0.5em}
\textbf{Wrong:}
\begin{itemize}
\item B: does not list all packages.
\item C: not a dependency file.
\item D: not generated.
\end{itemize}
\end{frame}

% ---------------------------------------------------------------
\begin{frame}{Question 28}
How do you implement a custom model with preprocessing and postprocessing logic in MLflow?
\begin{enumerate}[label=\Alph*.]
\item Subclass \texttt{mlflow.pyfunc.PythonModel} and implement \texttt{predict()}
\item Subclass \texttt{mlflow.Run}
\item Override \texttt{mlflow.start\_run}
\item Write a Delta trigger
\end{enumerate}
\end{frame}
\begin{frame}{Answer 28}
\textbf{Correct: A} --- a custom pyfunc wraps arbitrary logic.\\
\vspace{0.5em}
\textbf{Wrong:}
\begin{itemize}
\item B: not for models.
\item C: not for models.
\item D: not for models.
\end{itemize}
\end{frame}

% ---------------------------------------------------------------
\begin{frame}{Question 29}
Which call evaluates a logged model with default metrics and plots?
\begin{enumerate}[label=\Alph*.]
\item \texttt{mlflow.validate\_run()}
\item \texttt{mlflow.score()}
\item \texttt{mlflow.test\_model()}
\item \texttt{mlflow.evaluate(...)}
\end{enumerate}
\end{frame}
\begin{frame}{Answer 29}
\textbf{Correct: D} --- it computes metrics and artifacts for a model and dataset.\\
\vspace{0.5em}
\textbf{Wrong:}
\begin{itemize}
\item A: not an API.
\item B: not an API.
\item C: not an API.
\end{itemize}
\end{frame}

% ---------------------------------------------------------------
\begin{frame}{Question 30}
How can you disable MLflow autologging?
\begin{enumerate}[label=\Alph*.]
\item \texttt{mlflow.stop\_autolog()}
\item \texttt{mlflow.autolog(disable=True)}
\item \texttt{mlflow.end\_run(autolog=False)}
\item Delete the experiment
\end{enumerate}
\end{frame}
\begin{frame}{Answer 30}
\textbf{Correct: B} --- the disable flag turns autologging off.\\
\vspace{0.5em}
\textbf{Wrong:}
\begin{itemize}
\item A: not an API.
\item C: not a parameter.
\item D: not needed.
\end{itemize}
\end{frame}

% ---------------------------------------------------------------
\begin{frame}{Question 31}
Where does MLflow Tracking on Databricks store run metadata by default?
\begin{enumerate}[label=\Alph*.]
\item In a Git repository
\item Only on the cluster's local disk
\item In the feature store
\item In the workspace-managed tracking server for the experiment
\end{enumerate}
\end{frame}
\begin{frame}{Answer 31}
\textbf{Correct: D} --- Databricks hosts tracking for you.\\
\vspace{0.5em}
\textbf{Wrong:}
\begin{itemize}
\item A: not used for runs.
\item B: would be lost on termination.
\item C: not used for runs.
\end{itemize}
\end{frame}

% ---------------------------------------------------------------
\begin{frame}{Question 32}
Which client method lists versions of a registered model?
\begin{enumerate}[label=\Alph*.]
\item \texttt{MlflowClient().search\_model\_versions("name='catalog.schema.model'")}
\item \texttt{mlflow.search\_runs()}
\item \texttt{MlflowClient().list\_versions\_df()}
\item \texttt{mlflow.get\_versions()}
\end{enumerate}
\end{frame}
\begin{frame}{Answer 32}
\textbf{Correct: A} --- it filters model versions by name.\\
\vspace{0.5em}
\textbf{Wrong:}
\begin{itemize}
\item B: searches runs.
\item C: not an API.
\item D: not an API.
\end{itemize}
\end{frame}

% ---------------------------------------------------------------
\begin{frame}{Question 33}
Which tag or name argument sets a friendly name for a run?
\begin{enumerate}[label=\Alph*.]
\item \texttt{mlflow.end\_run("baseline")}
\item \texttt{mlflow.start\_run(label="baseline")}
\item \texttt{mlflow.start\_run(run\_name="baseline")}
\item \texttt{mlflow.log\_param("name")}
\end{enumerate}
\end{frame}
\begin{frame}{Answer 33}
\textbf{Correct: C} --- run\_name labels the run.\\
\vspace{0.5em}
\textbf{Wrong:}
\begin{itemize}
\item A: ends a run.
\item B: not an argument.
\item D: logs a param.
\end{itemize}
\end{frame}

% ---------------------------------------------------------------
\begin{frame}{Question 34}
What does \texttt{mlflow.end\_run()} do?
\begin{enumerate}[label=\Alph*.]
\item Deletes the experiment
\item Stops the cluster
\item Ends the active run
\item Registers the model
\end{enumerate}
\end{frame}
\begin{frame}{Answer 34}
\textbf{Correct: C} --- a run ends when the context exits or when end\_run is called.\\
\vspace{0.5em}
\textbf{Wrong:}
\begin{itemize}
\item A: not what it does.
\item B: not what it does.
\item D: not what it does.
\end{itemize}
\end{frame}

% ---------------------------------------------------------------
\begin{frame}{Question 35}
In Unity Catalog, what must exist before a model can be registered?
\begin{enumerate}[label=\Alph*.]
\item A serving endpoint
\item A feature table
\item A catalog and schema where the user has CREATE MODEL
\item An AutoML experiment
\end{enumerate}
\end{frame}
\begin{frame}{Answer 35}
\textbf{Correct: C} --- models live inside schemas.\\
\vspace{0.5em}
\textbf{Wrong:}
\begin{itemize}
\item A: created later.
\item B: not required.
\item D: not required.
\end{itemize}
\end{frame}

% ---------------------------------------------------------------
\begin{frame}{Question 36}
How can you attach metadata such as validation status to a model version?
\begin{enumerate}[label=\Alph*.]
\item Rename the experiment
\item Change the Delta schema
\item Use model version tags, for example \texttt{set\_model\_version\_tag}
\item Edit the cluster policy
\end{enumerate}
\end{frame}
\begin{frame}{Answer 36}
\textbf{Correct: C} --- tags store key-value metadata.\\
\vspace{0.5em}
\textbf{Wrong:}
\begin{itemize}
\item A: does not tag the version.
\item B: unrelated.
\item D: unrelated.
\end{itemize}
\end{frame}

% ---------------------------------------------------------------
\begin{frame}{Question 37}
What is a common practice before moving the Champion alias to a new version?
\begin{enumerate}[label=\Alph*.]
\item Validate the new version against checks and a baseline
\item Skip tests for speed
\item Delete the old version immediately
\item Retrain without logging
\end{enumerate}
\end{frame}
\begin{frame}{Answer 37}
\textbf{Correct: A} --- validation reduces production risk.\\
\vspace{0.5em}
\textbf{Wrong:}
\begin{itemize}
\item B: risky.
\item C: prevents rollback.
\item D: loses traceability.
\end{itemize}
\end{frame}

% ---------------------------------------------------------------
\begin{frame}{Question 38}
Which URI loads a specific numbered version of a UC registered model?
\begin{enumerate}[label=\Alph*.]
\item \texttt{runs:/model/3}
\item \texttt{models:/3/catalog.schema.model}
\item \texttt{versions:/model@3}
\item \texttt{models:/catalog.schema.model/3}
\end{enumerate}
\end{frame}
\begin{frame}{Answer 38}
\textbf{Correct: D} --- the version number follows the model name.\\
\vspace{0.5em}
\textbf{Wrong:}
\begin{itemize}
\item A: not a model version URI.
\item B: wrong order.
\item C: not valid.
\end{itemize}
\end{frame}

% ---------------------------------------------------------------
\begin{frame}{Question 39}
Which statement about aliases is true?
\begin{enumerate}[label=\Alph*.]
\item One alias points to one version at a time within a model
\item An alias can point to many versions at once
\item Aliases are only for the legacy registry
\item Aliases delete old versions
\end{enumerate}
\end{frame}
\begin{frame}{Answer 39}
\textbf{Correct: A} --- reassigning moves the alias.\\
\vspace{0.5em}
\textbf{Wrong:}
\begin{itemize}
\item B: it points to one.
\item C: they are the UC mechanism.
\item D: they do not.
\end{itemize}
\end{frame}

% ---------------------------------------------------------------
\begin{frame}{Question 40}
A model in Unity Catalog was trained on features from a UC feature table. What does UC lineage show?
\begin{enumerate}[label=\Alph*.]
\item The upstream feature table and data sources linked to the model
\item Only the run duration
\item Nothing, models have no lineage
\item Only the cluster size
\end{enumerate}
\end{frame}
\begin{frame}{Answer 40}
\textbf{Correct: A} --- lineage connects data assets and models.\\
\vspace{0.5em}
\textbf{Wrong:}
\begin{itemize}
\item B: not lineage.
\item C: UC tracks model lineage.
\item D: not lineage.
\end{itemize}
\end{frame}

% ---------------------------------------------------------------
\begin{frame}{Question 41}
Which call reads a feature table as a Spark DataFrame?
\begin{enumerate}[label=\Alph*.]
\item \texttt{fe.read\_table(name="catalog.schema.features")}
\item \texttt{fe.create\_table()}
\item \texttt{fe.load\_features()}
\item \texttt{fe.drop\_table()}
\end{enumerate}
\end{frame}
\begin{frame}{Answer 41}
\textbf{Correct: A} --- read\_table returns the table data.\\
\vspace{0.5em}
\textbf{Wrong:}
\begin{itemize}
\item B: creates a table.
\item C: not an API.
\item D: deletes a table.
\end{itemize}
\end{frame}

% ---------------------------------------------------------------
\begin{frame}{Question 42}
Which call deletes a feature table?
\begin{enumerate}[label=\Alph*.]
\item \texttt{fe.write\_table(mode="delete")}
\item \texttt{fe.remove()}
\item \texttt{fe.drop\_table(name=...)}
\item \texttt{fe.delete\_features()}
\end{enumerate}
\end{frame}
\begin{frame}{Answer 42}
\textbf{Correct: C} --- drop\_table removes the feature table.\\
\vspace{0.5em}
\textbf{Wrong:}
\begin{itemize}
\item A: not a mode.
\item B: not an API.
\item D: not an API.
\end{itemize}
\end{frame}

% ---------------------------------------------------------------
\begin{frame}{Question 43}
What does \texttt{exclude\_columns} do in \texttt{create\_training\_set}?
\begin{enumerate}[label=\Alph*.]
\item Excludes tables from lineage
\item Excludes failed runs
\item Removes listed columns from the resulting training DataFrame
\item Excludes rows with nulls
\end{enumerate}
\end{frame}
\begin{frame}{Answer 43}
\textbf{Correct: C} --- use it to drop keys or other non-feature columns.\\
\vspace{0.5em}
\textbf{Wrong:}
\begin{itemize}
\item A: unrelated.
\item B: unrelated.
\item D: it drops columns.
\end{itemize}
\end{frame}

% ---------------------------------------------------------------
\begin{frame}{Question 44}
Which problem does a feature store help reduce?
\begin{enumerate}[label=\Alph*.]
\item Cluster costs only
\item Label noise
\item Training-serving skew caused by inconsistent feature code
\item Model overfitting
\end{enumerate}
\end{frame}
\begin{frame}{Answer 44}
\textbf{Correct: C} --- one definition of features is used everywhere.\\
\vspace{0.5em}
\textbf{Wrong:}
\begin{itemize}
\item A: not its purpose.
\item B: not directly solved.
\item D: not directly solved.
\end{itemize}
\end{frame}

% ---------------------------------------------------------------
\begin{frame}{Question 45}
A feature is computed from request-time input such as the current transaction amount. What helps?
\begin{enumerate}[label=\Alph*.]
\item An on-demand feature function computed at inference
\item Dropping the feature
\item Storing it in a notebook cell
\item A nightly batch feature only
\end{enumerate}
\end{frame}
\begin{frame}{Answer 45}
\textbf{Correct: A} --- feature functions compute values that cannot be precomputed.\\
\vspace{0.5em}
\textbf{Wrong:}
\begin{itemize}
\item B: loses information.
\item C: not governed.
\item D: would be stale.
\end{itemize}
\end{frame}

% ---------------------------------------------------------------
\begin{frame}{Question 46}
What is a feature table in Unity Catalog, technically?
\begin{enumerate}[label=\Alph*.]
\item A serving endpoint
\item A CSV file in DBFS
\item A pandas DataFrame in memory
\item A Delta table with a primary key constraint
\end{enumerate}
\end{frame}
\begin{frame}{Answer 46}
\textbf{Correct: D} --- UC feature tables are Delta tables with keys.\\
\vspace{0.5em}
\textbf{Wrong:}
\begin{itemize}
\item A: not a table.
\item B: not a feature table.
\item C: not persisted.
\end{itemize}
\end{frame}

% ---------------------------------------------------------------
\begin{frame}{Question 47}
Which parameter provides the join key column when defining a \texttt{FeatureLookup}?
\begin{enumerate}[label=\Alph*.]
\item \texttt{lookup\_key}
\item \texttt{join\_col\_name}
\item \texttt{label}
\item \texttt{primary\_keys}
\end{enumerate}
\end{frame}
\begin{frame}{Answer 47}
\textbf{Correct: A} --- it names the key column in the label DataFrame.\\
\vspace{0.5em}
\textbf{Wrong:}
\begin{itemize}
\item B: not the parameter.
\item C: used in the training set.
\item D: used when creating tables.
\end{itemize}
\end{frame}

% ---------------------------------------------------------------
\begin{frame}{Question 48}
Which \texttt{FeatureLookup} parameter enables point-in-time joins?
\begin{enumerate}[label=\Alph*.]
\item \texttt{feature\_names}
\item \texttt{timestamp\_lookup\_key}
\item \texttt{version\_as\_of}
\item \texttt{time\_travel=True}
\end{enumerate}
\end{frame}
\begin{frame}{Answer 48}
\textbf{Correct: B} --- it names the timestamp column in the label data.\\
\vspace{0.5em}
\textbf{Wrong:}
\begin{itemize}
\item A: lists features only.
\item C: not a parameter.
\item D: not a parameter.
\end{itemize}
\end{frame}

% ---------------------------------------------------------------
\begin{frame}{Question 49}
What is the difference between normalization (min-max) and standardization?
\begin{enumerate}[label=\Alph*.]
\item Min-max rescales to a fixed range; standardization centers to mean 0 and unit variance
\item They are identical
\item Standardization maps values to 0 or 1
\item Min-max centers at the median
\end{enumerate}
\end{frame}
\begin{frame}{Answer 49}
\textbf{Correct: A} --- different scaling approaches for different needs.\\
\vspace{0.5em}
\textbf{Wrong:}
\begin{itemize}
\item B: they differ.
\item C: that is binarization.
\item D: not true.
\end{itemize}
\end{frame}

% ---------------------------------------------------------------
\begin{frame}{Question 50}
Which Spark ML class standardizes features?
\begin{enumerate}[label=\Alph*.]
\item \texttt{StandardScaler}
\item \texttt{VectorAssembler}
\item \texttt{StringIndexer}
\item \texttt{Tokenizer}
\end{enumerate}
\end{frame}
\begin{frame}{Answer 50}
\textbf{Correct: A} --- it scales to unit variance and optionally centers.\\
\vspace{0.5em}
\textbf{Wrong:}
\begin{itemize}
\item B: combines columns.
\item C: indexes categories.
\item D: splits text.
\end{itemize}
\end{frame}

% ---------------------------------------------------------------
\begin{frame}{Question 51}
Which Spark ML class rescales features to a given range such as 0 to 1?
\begin{enumerate}[label=\Alph*.]
\item \texttt{Bucketizer}
\item \texttt{PCA}
\item \texttt{Imputer}
\item \texttt{MinMaxScaler}
\end{enumerate}
\end{frame}
\begin{frame}{Answer 51}
\textbf{Correct: D} --- it does min-max scaling.\\
\vspace{0.5em}
\textbf{Wrong:}
\begin{itemize}
\item A: bins continuous values.
\item B: reduces dimensions.
\item C: fills missing values.
\end{itemize}
\end{frame}

% ---------------------------------------------------------------
\begin{frame}{Question 52}
Which Spark ML class bins a continuous column using specified split points?
\begin{enumerate}[label=\Alph*.]
\item \texttt{StringIndexer}
\item \texttt{OneHotEncoder}
\item \texttt{StandardScaler}
\item \texttt{Bucketizer}
\end{enumerate}
\end{frame}
\begin{frame}{Answer 52}
\textbf{Correct: D} --- it maps values into buckets.\\
\vspace{0.5em}
\textbf{Wrong:}
\begin{itemize}
\item A: indexes strings.
\item B: encodes indexes.
\item C: scales values.
\end{itemize}
\end{frame}

% ---------------------------------------------------------------
\begin{frame}{Question 53}
What does PCA do?
\begin{enumerate}[label=\Alph*.]
\item Balances classes
\item Projects data to fewer uncorrelated components that capture most variance
\item Tunes hyperparameters
\item Fills missing values
\end{enumerate}
\end{frame}
\begin{frame}{Answer 53}
\textbf{Correct: B} --- it reduces dimensionality.\\
\vspace{0.5em}
\textbf{Wrong:}
\begin{itemize}
\item A: not its purpose.
\item C: not its purpose.
\item D: not its purpose.
\end{itemize}
\end{frame}

% ---------------------------------------------------------------
\begin{frame}{Question 54}
What problem does multicollinearity cause for a linear model?
\begin{enumerate}[label=\Alph*.]
\item It improves interpretability
\item It removes the intercept
\item It makes training impossible on Spark
\item Unstable and hard-to-interpret coefficients
\end{enumerate}
\end{frame}
\begin{frame}{Answer 54}
\textbf{Correct: D} --- correlated predictors make coefficients unreliable.\\
\vspace{0.5em}
\textbf{Wrong:}
\begin{itemize}
\item A: it harms it.
\item B: not true.
\item C: not true.
\end{itemize}
\end{frame}

% ---------------------------------------------------------------
\begin{frame}{Question 55}
Which measure is commonly used to detect multicollinearity?
\begin{enumerate}[label=\Alph*.]
\item ROC AUC
\item Silhouette score
\item Log loss
\item Variance inflation factor (VIF)
\end{enumerate}
\end{frame}
\begin{frame}{Answer 55}
\textbf{Correct: D} --- a high VIF suggests collinearity.\\
\vspace{0.5em}
\textbf{Wrong:}
\begin{itemize}
\item A: classification metric.
\item B: clustering metric.
\item C: classification loss.
\end{itemize}
\end{frame}

% ---------------------------------------------------------------
\begin{frame}{Question 56}
A nominal categorical feature such as color has no natural order. Which encoding is appropriate for a linear model?
\begin{enumerate}[label=\Alph*.]
\item Dropping the column always
\item Squaring the values
\item Ordinal integer encoding
\item One-hot encoding
\end{enumerate}
\end{frame}
\begin{frame}{Answer 56}
\textbf{Correct: D} --- it avoids imposing a false order.\\
\vspace{0.5em}
\textbf{Wrong:}
\begin{itemize}
\item A: loses information.
\item B: not an encoding.
\item C: implies a false order.
\end{itemize}
\end{frame}

% ---------------------------------------------------------------
\begin{frame}{Question 57}
Why can target encoding cause leakage?
\begin{enumerate}[label=\Alph*.]
\item It always lowers accuracy
\item It removes all labels
\item It requires a GPU
\item It uses label information that can leak into training features if not computed carefully
\end{enumerate}
\end{frame}
\begin{frame}{Answer 57}
\textbf{Correct: D} --- compute it within folds or on training data only.\\
\vspace{0.5em}
\textbf{Wrong:}
\begin{itemize}
\item A: not true.
\item B: not true.
\item C: not true.
\end{itemize}
\end{frame}

% ---------------------------------------------------------------
\begin{frame}{Question 58}
A feature has 40 percent missing values. Which step is reasonable?
\begin{enumerate}[label=\Alph*.]
\item Investigate why, then impute or add a missing indicator, or drop if uninformative
\item Ignore it silently
\item Always drop all such rows
\item Always fill with zero
\end{enumerate}
\end{frame}
\begin{frame}{Answer 58}
\textbf{Correct: A} --- choose based on cause and importance.\\
\vspace{0.5em}
\textbf{Wrong:}
\begin{itemize}
\item B: can break models.
\item C: may lose too much data.
\item D: can distort the data.
\end{itemize}
\end{frame}

% ---------------------------------------------------------------
\begin{frame}{Question 59}
Which operation removes duplicate rows in Spark?
\begin{enumerate}[label=\Alph*.]
\item \texttt{df.dropDuplicates()}
\item \texttt{df.dedupe()}
\item \texttt{df.removeRows()}
\item \texttt{df.unique()}
\end{enumerate}
\end{frame}
\begin{frame}{Answer 59}
\textbf{Correct: A} --- it removes duplicates, optionally on a subset of columns.\\
\vspace{0.5em}
\textbf{Wrong:}
\begin{itemize}
\item B: not an API.
\item C: not an API.
\item D: not an API.
\end{itemize}
\end{frame}

% ---------------------------------------------------------------
\begin{frame}{Question 60}
Which Spark call removes rows containing nulls?
\begin{enumerate}[label=\Alph*.]
\item \texttt{df.fillna()}
\item \texttt{df.dropDuplicates()}
\item \texttt{df.dropna()} (\texttt{df.na.drop()})
\item \texttt{df.nullify()}
\end{enumerate}
\end{frame}
\begin{frame}{Answer 60}
\textbf{Correct: C} --- it drops rows with missing values.\\
\vspace{0.5em}
\textbf{Wrong:}
\begin{itemize}
\item A: fills nulls.
\item B: removes duplicates.
\item D: not an API.
\end{itemize}
\end{frame}

% ---------------------------------------------------------------
\begin{frame}{Question 61}
Which call fills nulls with a constant in a Spark DataFrame?
\begin{enumerate}[label=\Alph*.]
\item \texttt{df.fillna(0)} (\texttt{df.na.fill(0)})
\item \texttt{df.replaceNull(0)}
\item \texttt{df.dropna(0)}
\item \texttt{df.coalesce(0)}
\end{enumerate}
\end{frame}
\begin{frame}{Answer 61}
\textbf{Correct: A} --- it replaces nulls with the given value.\\
\vspace{0.5em}
\textbf{Wrong:}
\begin{itemize}
\item B: not an API.
\item C: drops rows.
\item D: changes partitions.
\end{itemize}
\end{frame}

% ---------------------------------------------------------------
\begin{frame}{Question 62}
Which pandas UDF type lets you apply Python logic to batches of a column?
\begin{enumerate}[label=\Alph*.]
\item A Delta trigger
\item A Series-to-Series pandas UDF
\item A row-at-a-time Python UDF
\item A SQL view
\end{enumerate}
\end{frame}
\begin{frame}{Answer 62}
\textbf{Correct: B} --- it receives and returns a pandas Series per batch.\\
\vspace{0.5em}
\textbf{Wrong:}
\begin{itemize}
\item A: not a UDF.
\item C: is not vectorized.
\item D: not a UDF.
\end{itemize}
\end{frame}

% ---------------------------------------------------------------
\begin{frame}{Question 63}
What should a Hyperopt objective function return so \texttt{fmin} can record status and loss?
\begin{enumerate}[label=\Alph*.]
\item A list of models
\item Nothing
\item A dict such as \texttt{\{"loss": value, "status": STATUS\_OK\}}
\item A pandas DataFrame
\end{enumerate}
\end{frame}
\begin{frame}{Answer 63}
\textbf{Correct: C} --- fmin reads the loss and status keys.\\
\vspace{0.5em}
\textbf{Wrong:}
\begin{itemize}
\item A: not accepted.
\item B: a loss is required.
\item D: not accepted.
\end{itemize}
\end{frame}

% ---------------------------------------------------------------
\begin{frame}{Question 64}
In \texttt{hp.loguniform("lr", low, high)}, what are \texttt{low} and \texttt{high}?
\begin{enumerate}[label=\Alph*.]
\item The mean and std
\item The natural logs of the range limits
\item The actual range limits
\item The number of trials
\end{enumerate}
\end{frame}
\begin{frame}{Answer 64}
\textbf{Correct: B} --- so use np.log(1e-5) and np.log(1e-1) for 1e-5 to 1e-1.\\
\vspace{0.5em}
\textbf{Wrong:}
\begin{itemize}
\item A: not their meaning.
\item C: they are logs.
\item D: not their meaning.
\end{itemize}
\end{frame}

% ---------------------------------------------------------------
\begin{frame}{Question 65}
What does \texttt{max\_evals} control in \texttt{fmin}?
\begin{enumerate}[label=\Alph*.]
\item The number of CV folds
\item The number of workers
\item The tree depth
\item The maximum number of trials to evaluate
\end{enumerate}
\end{frame}
\begin{frame}{Answer 65}
\textbf{Correct: D} --- it bounds the search budget.\\
\vspace{0.5em}
\textbf{Wrong:}
\begin{itemize}
\item A: set in the objective.
\item B: that is parallelism.
\item C: a model hyperparameter.
\end{itemize}
\end{frame}

% ---------------------------------------------------------------
\begin{frame}{Question 66}
Which function converts the index results from \texttt{fmin} back to the actual parameter values?
\begin{enumerate}[label=\Alph*.]
\item \texttt{hp.resolve()}
\item \texttt{space\_eval(space, best)}
\item \texttt{tpe.decode()}
\item \texttt{eval\_space()}
\end{enumerate}
\end{frame}
\begin{frame}{Answer 66}
\textbf{Correct: B} --- it maps hp.choice indexes to values.\\
\vspace{0.5em}
\textbf{Wrong:}
\begin{itemize}
\item A: not an API.
\item C: not an API.
\item D: not an API.
\end{itemize}
\end{frame}

% ---------------------------------------------------------------
\begin{frame}{Question 67}
Which Hyperopt object stores all trial results for later inspection?
\begin{enumerate}[label=\Alph*.]
\item \texttt{STATUS\_OK}
\item \texttt{Trials} (or \texttt{SparkTrials})
\item \texttt{tpe.suggest}
\item \texttt{hp.choice}
\end{enumerate}
\end{frame}
\begin{frame}{Answer 67}
\textbf{Correct: B} --- pass it to fmin to keep trial info.\\
\vspace{0.5em}
\textbf{Wrong:}
\begin{itemize}
\item A: a status constant.
\item C: an algorithm.
\item D: defines a space.
\end{itemize}
\end{frame}

% ---------------------------------------------------------------
\begin{frame}{Question 68}
Why should the objective function evaluate on validation data, not training data?
\begin{enumerate}[label=\Alph*.]
\item Training loss rewards overfit hyperparameters
\item Training data is too small
\item Hyperopt forbids training data
\item Validation data is faster
\end{enumerate}
\end{frame}
\begin{frame}{Answer 68}
\textbf{Correct: A} --- validation estimates generalization.\\
\vspace{0.5em}
\textbf{Wrong:}
\begin{itemize}
\item B: not the reason.
\item C: not a rule.
\item D: not the reason.
\end{itemize}
\end{frame}

% ---------------------------------------------------------------
\begin{frame}{Question 69}
What does \texttt{numFolds} set in Spark ML \texttt{CrossValidator}?
\begin{enumerate}[label=\Alph*.]
\item The number of workers
\item The number of folds
\item The number of features
\item The number of trees
\end{enumerate}
\end{frame}
\begin{frame}{Answer 69}
\textbf{Correct: B} --- the data is split into that many folds.\\
\vspace{0.5em}
\textbf{Wrong:}
\begin{itemize}
\item A: not set here.
\item C: not set here.
\item D: a model parameter.
\end{itemize}
\end{frame}

% ---------------------------------------------------------------
\begin{frame}{Question 70}
How do you build a parameter grid in Spark ML?
\begin{enumerate}[label=\Alph*.]
\item \texttt{GridSearch.build()}
\item \texttt{ParamGridBuilder().addGrid(param, values).build()}
\item \texttt{Hyperopt.grid()}
\item \texttt{CrossValidator.grid()}
\end{enumerate}
\end{frame}
\begin{frame}{Answer 70}
\textbf{Correct: B} --- the builder creates the combinations.\\
\vspace{0.5em}
\textbf{Wrong:}
\begin{itemize}
\item A: not Spark ML.
\item C: not an API.
\item D: not an API.
\end{itemize}
\end{frame}

% ---------------------------------------------------------------
\begin{frame}{Question 71}
After fitting a \texttt{CrossValidator}, how do you access the best model?
\begin{enumerate}[label=\Alph*.]
\item \texttt{cv\_model.best\_run}
\item \texttt{cv\_model.top}
\item \texttt{cv\_model.bestModel}
\item \texttt{cv\_model.winner}
\end{enumerate}
\end{frame}
\begin{frame}{Answer 71}
\textbf{Correct: C} --- the fitted CrossValidatorModel exposes bestModel.\\
\vspace{0.5em}
\textbf{Wrong:}
\begin{itemize}
\item A: not an attribute.
\item B: not an attribute.
\item D: not an attribute.
\end{itemize}
\end{frame}

% ---------------------------------------------------------------
\begin{frame}{Question 72}
What does \texttt{avgMetrics} contain in a \texttt{CrossValidatorModel}?
\begin{enumerate}[label=\Alph*.]
\item The feature importances
\item The list of columns
\item The average cross-validation metric for each parameter combination
\item The training time per fold
\end{enumerate}
\end{frame}
\begin{frame}{Answer 72}
\textbf{Correct: C} --- use it to compare grid points.\\
\vspace{0.5em}
\textbf{Wrong:}
\begin{itemize}
\item A: not contained.
\item B: not contained.
\item D: not contained.
\end{itemize}
\end{frame}

% ---------------------------------------------------------------
\begin{frame}{Question 73}
Which Spark ML class implements logistic regression?
\begin{enumerate}[label=\Alph*.]
\item \texttt{LinearRegression}
\item \texttt{LogisticRegression}
\item \texttt{PCA}
\item \texttt{KMeans}
\end{enumerate}
\end{frame}
\begin{frame}{Answer 73}
\textbf{Correct: B} --- it trains binary and multiclass classifiers.\\
\vspace{0.5em}
\textbf{Wrong:}
\begin{itemize}
\item A: predicts continuous targets.
\item C: reduces dimensions.
\item D: is clustering.
\end{itemize}
\end{frame}

% ---------------------------------------------------------------
\begin{frame}{Question 74}
In Spark ML regularization, what does \texttt{elasticNetParam=1.0} mean?
\begin{enumerate}[label=\Alph*.]
\item A 50/50 mix
\item No regularization
\item Pure L2 penalty
\item Pure L1 (Lasso) penalty
\end{enumerate}
\end{frame}
\begin{frame}{Answer 74}
\textbf{Correct: D} --- 0 is L2 and 1 is L1.\\
\vspace{0.5em}
\textbf{Wrong:}
\begin{itemize}
\item A: that is 0.5.
\item B: that is \texttt{regParam=0}.
\item C: that is 0.
\end{itemize}
\end{frame}

% ---------------------------------------------------------------
\begin{frame}{Question 75}
Which parameter sets the regularization strength in Spark ML linear models?
\begin{enumerate}[label=\Alph*.]
\item \texttt{featuresCol}
\item \texttt{numTrees}
\item \texttt{regParam}
\item \texttt{maxBins}
\end{enumerate}
\end{frame}
\begin{frame}{Answer 75}
\textbf{Correct: C} --- larger values mean stronger regularization.\\
\vspace{0.5em}
\textbf{Wrong:}
\begin{itemize}
\item A: column name.
\item B: forest size.
\item D: tree discretization.
\end{itemize}
\end{frame}

% ---------------------------------------------------------------
\begin{frame}{Question 76}
Which parameter controls maximum tree depth in Spark ML decision trees?
\begin{enumerate}[label=\Alph*.]
\item \texttt{maxDepth}
\item \texttt{maxIter}
\item \texttt{numFolds}
\item \texttt{regParam}
\end{enumerate}
\end{frame}
\begin{frame}{Answer 76}
\textbf{Correct: A} --- deeper trees can overfit.\\
\vspace{0.5em}
\textbf{Wrong:}
\begin{itemize}
\item B: iterations.
\item C: CV folds.
\item D: regularization.
\end{itemize}
\end{frame}

% ---------------------------------------------------------------
\begin{frame}{Question 77}
Which Spark ML estimator performs clustering?
\begin{enumerate}[label=\Alph*.]
\item \texttt{LogisticRegression}
\item \texttt{DecisionTreeClassifier}
\item \texttt{KMeans}
\item \texttt{GBTRegressor}
\end{enumerate}
\end{frame}
\begin{frame}{Answer 77}
\textbf{Correct: C} --- it groups rows into k clusters.\\
\vspace{0.5em}
\textbf{Wrong:}
\begin{itemize}
\item A: supervised.
\item B: supervised.
\item D: supervised.
\end{itemize}
\end{frame}

% ---------------------------------------------------------------
\begin{frame}{Question 78}
Which evaluator should you use to compute silhouette score for clustering in Spark ML?
\begin{enumerate}[label=\Alph*.]
\item \texttt{BinaryClassificationEvaluator}
\item \texttt{RegressionEvaluator}
\item \texttt{MulticlassClassificationEvaluator}
\item \texttt{ClusteringEvaluator}
\end{enumerate}
\end{frame}
\begin{frame}{Answer 78}
\textbf{Correct: D} --- its default metric is silhouette.\\
\vspace{0.5em}
\textbf{Wrong:}
\begin{itemize}
\item A: binary classification.
\item B: regression metrics.
\item C: classification metrics.
\end{itemize}
\end{frame}

% ---------------------------------------------------------------
\begin{frame}{Question 79}
Which Spark ML algorithm builds recommendations with collaborative filtering?
\begin{enumerate}[label=\Alph*.]
\item \texttt{PCA}
\item \texttt{ALS}
\item \texttt{LinearSVC}
\item \texttt{KMeans}
\end{enumerate}
\end{frame}
\begin{frame}{Answer 79}
\textbf{Correct: B} --- alternating least squares for user-item ratings.\\
\vspace{0.5em}
\textbf{Wrong:}
\begin{itemize}
\item A: dimensionality reduction.
\item C: classification.
\item D: clustering.
\end{itemize}
\end{frame}

% ---------------------------------------------------------------
\begin{frame}{Question 80}
Which statement about fitting preprocessing in a pipeline is correct?
\begin{enumerate}[label=\Alph*.]
\item Fit separately on each split
\item Fit on training data only and apply the fitted transformers to validation and test data
\item Fit on the full dataset first
\item Fit on test data only
\end{enumerate}
\end{frame}
\begin{frame}{Answer 80}
\textbf{Correct: B} --- this prevents leakage.\\
\vspace{0.5em}
\textbf{Wrong:}
\begin{itemize}
\item A: creates inconsistent features.
\item C: leaks information.
\item D: defeats the purpose.
\end{itemize}
\end{frame}

% ---------------------------------------------------------------
\begin{frame}{Question 81}
What is precision?
\begin{enumerate}[label=\Alph*.]
\item TP / (TP + FP)
\item TP / (TP + FN)
\item (TP + TN) / total
\item TN / (TN + FP)
\end{enumerate}
\end{frame}
\begin{frame}{Answer 81}
\textbf{Correct: A} --- the share of predicted positives that are correct.\\
\vspace{0.5em}
\textbf{Wrong:}
\begin{itemize}
\item B: that is recall.
\item C: that is accuracy.
\item D: that is specificity.
\end{itemize}
\end{frame}

% ---------------------------------------------------------------
\begin{frame}{Question 82}
What is recall?
\begin{enumerate}[label=\Alph*.]
\item TP / (TP + FP)
\item (TP + TN) / total
\item FP / (FP + TN)
\item TP / (TP + FN)
\end{enumerate}
\end{frame}
\begin{frame}{Answer 82}
\textbf{Correct: D} --- the share of actual positives that are found.\\
\vspace{0.5em}
\textbf{Wrong:}
\begin{itemize}
\item A: that is precision.
\item B: that is accuracy.
\item C: that is false positive rate.
\end{itemize}
\end{frame}

% ---------------------------------------------------------------
\begin{frame}{Question 83}
What does the F1 score combine?
\begin{enumerate}[label=\Alph*.]
\item The arithmetic mean of accuracy and AUC
\item The harmonic mean of precision and recall
\item The sum of TP and TN
\item The maximum of precision and recall
\end{enumerate}
\end{frame}
\begin{frame}{Answer 83}
\textbf{Correct: B} --- it balances both.\\
\vspace{0.5em}
\textbf{Wrong:}
\begin{itemize}
\item A: not its definition.
\item C: not its definition.
\item D: not its definition.
\end{itemize}
\end{frame}

% ---------------------------------------------------------------
\begin{frame}{Question 84}
Which is true about MSE versus MAE units?
\begin{enumerate}[label=\Alph*.]
\item MAE is squared
\item Both are unitless
\item MSE is in squared units; MAE is in the target's units
\item MSE equals RMSE
\end{enumerate}
\end{frame}
\begin{frame}{Answer 84}
\textbf{Correct: C} --- RMSE returns to the original units.\\
\vspace{0.5em}
\textbf{Wrong:}
\begin{itemize}
\item A: MSE is squared.
\item B: not true.
\item D: RMSE is the square root of MSE.
\end{itemize}
\end{frame}

% ---------------------------------------------------------------
\begin{frame}{Question 85}
Why can MAPE be a poor metric?
\begin{enumerate}[label=\Alph*.]
\item It ignores errors
\item It cannot be computed in Spark
\item It is undefined or unstable when actual values are zero or near zero
\item It only works for classification
\end{enumerate}
\end{frame}
\begin{frame}{Answer 85}
\textbf{Correct: C} --- it divides by the actual value.\\
\vspace{0.5em}
\textbf{Wrong:}
\begin{itemize}
\item A: not true.
\item B: not the reason.
\item D: it is for regression.
\end{itemize}
\end{frame}

% ---------------------------------------------------------------
\begin{frame}{Question 86}
What does adjusted R2 do that R2 does not?
\begin{enumerate}[label=\Alph*.]
\item Predicts probabilities
\item Penalizes adding predictors that do not improve the model
\item Ignores the number of features
\item Replaces the test set
\end{enumerate}
\end{frame}
\begin{frame}{Answer 86}
\textbf{Correct: B} --- plain R2 never decreases when features are added.\\
\vspace{0.5em}
\textbf{Wrong:}
\begin{itemize}
\item A: unrelated.
\item C: R2 does that.
\item D: unrelated.
\end{itemize}
\end{frame}

% ---------------------------------------------------------------
\begin{frame}{Question 87}
What does log loss penalize heavily?
\begin{enumerate}[label=\Alph*.]
\item Missing values
\item Large feature counts
\item Confident but wrong probability predictions
\item Slow training
\end{enumerate}
\end{frame}
\begin{frame}{Answer 87}
\textbf{Correct: C} --- it uses the predicted probability of the true class.\\
\vspace{0.5em}
\textbf{Wrong:}
\begin{itemize}
\item A: not a loss factor.
\item B: not a loss factor.
\item D: not a loss factor.
\end{itemize}
\end{frame}

% ---------------------------------------------------------------
\begin{frame}{Question 88}
What is the difference between a parameter and a hyperparameter?
\begin{enumerate}[label=\Alph*.]
\item Hyperparameters are learned by gradient descent
\item They are the same
\item Parameters are set by the user
\item Parameters are learned from data; hyperparameters are set before training
\end{enumerate}
\end{frame}
\begin{frame}{Answer 88}
\textbf{Correct: D} --- for example coefficients versus tree depth.\\
\vspace{0.5em}
\textbf{Wrong:}
\begin{itemize}
\item A: they are set beforehand.
\item B: they differ.
\item C: that is hyperparameters.
\end{itemize}
\end{frame}

% ---------------------------------------------------------------
\begin{frame}{Question 89}
What are the roles of validation and test sets?
\begin{enumerate}[label=\Alph*.]
\item Validation guides tuning; test gives a final unbiased estimate
\item They are always identical
\item Test guides tuning; validation is final
\item Both are used for training
\end{enumerate}
\end{frame}
\begin{frame}{Answer 89}
\textbf{Correct: A} --- do not tune on the test set.\\
\vspace{0.5em}
\textbf{Wrong:}
\begin{itemize}
\item B: they must be separate.
\item C: reversed.
\item D: they are held out.
\end{itemize}
\end{frame}

% ---------------------------------------------------------------
\begin{frame}{Question 90}
What is the curse of dimensionality?
\begin{enumerate}[label=\Alph*.]
\item Training always becomes faster
\item Models always improve
\item Data becomes sparse and models harder to fit as the number of features grows
\item Features become perfectly correlated
\end{enumerate}
\end{frame}
\begin{frame}{Answer 90}
\textbf{Correct: C} --- more features need far more data.\\
\vspace{0.5em}
\textbf{Wrong:}
\begin{itemize}
\item A: not true.
\item B: not true.
\item D: not true.
\end{itemize}
\end{frame}

% ---------------------------------------------------------------
\begin{frame}{Question 91}
What is the effect of a gradient descent learning rate that is too large?
\begin{enumerate}[label=\Alph*.]
\item Training always converges faster and better
\item Nothing changes
\item The model underfits by design
\item Training can overshoot and diverge
\end{enumerate}
\end{frame}
\begin{frame}{Answer 91}
\textbf{Correct: D} --- steps are too big to converge.\\
\vspace{0.5em}
\textbf{Wrong:}
\begin{itemize}
\item A: not guaranteed.
\item B: the rate matters.
\item C: not the typical result.
\end{itemize}
\end{frame}

% ---------------------------------------------------------------
\begin{frame}{Question 92}
What does early stopping do?
\begin{enumerate}[label=\Alph*.]
\item Stops registering models
\item Stops training when validation performance stops improving
\item Stops reading data
\item Stops the cluster when idle
\end{enumerate}
\end{frame}
\begin{frame}{Answer 92}
\textbf{Correct: B} --- it helps prevent overfitting.\\
\vspace{0.5em}
\textbf{Wrong:}
\begin{itemize}
\item A: not related.
\item C: not related.
\item D: that is auto-termination.
\end{itemize}
\end{frame}

% ---------------------------------------------------------------
\begin{frame}{Question 93}
What happens with scale-to-zero on a Model Serving endpoint?
\begin{enumerate}[label=\Alph*.]
\item The registry is cleared
\item Predictions are cached forever
\item Compute scales down when there is no traffic, and the next request may see a cold start
\item The model is deleted
\end{enumerate}
\end{frame}
\begin{frame}{Answer 93}
\textbf{Correct: C} --- it saves cost at the price of latency.\\
\vspace{0.5em}
\textbf{Wrong:}
\begin{itemize}
\item A: not true.
\item B: not true.
\item D: the model remains.
\end{itemize}
\end{frame}

% ---------------------------------------------------------------
\begin{frame}{Question 94}
Which call queries a Databricks serving endpoint from Python using MLflow?
\begin{enumerate}[label=\Alph*.]
\item \texttt{mlflow.query\_endpoint()}
\item \texttt{mlflow.deployments.get\_deploy\_client("databricks").predict(endpoint=..., inputs=...)}
\item \texttt{spark.predict(endpoint)}
\item \texttt{mlflow.serve(endpoint)}
\end{enumerate}
\end{frame}
\begin{frame}{Answer 94}
\textbf{Correct: B} --- the deployments client calls the endpoint.\\
\vspace{0.5em}
\textbf{Wrong:}
\begin{itemize}
\item A: not an API.
\item C: not an API.
\item D: not valid.
\end{itemize}
\end{frame}

% ---------------------------------------------------------------
\begin{frame}{Question 95}
What authentication does a REST call to a Databricks serving endpoint need?
\begin{enumerate}[label=\Alph*.]
\item A bearer token (personal access token or OAuth token)
\item A Git token
\item No authentication
\item A cluster ID
\end{enumerate}
\end{frame}
\begin{frame}{Answer 95}
\textbf{Correct: A} --- the request carries an Authorization header.\\
\vspace{0.5em}
\textbf{Wrong:}
\begin{itemize}
\item B: not used.
\item C: endpoints are secured.
\item D: not used.
\end{itemize}
\end{frame}

% ---------------------------------------------------------------
\begin{frame}{Question 96}
In a multi-task Databricks job, how do you make task B run only after task A?
\begin{enumerate}[label=\Alph*.]
\item Put them in one cell
\item Use two clusters
\item Set task A as a dependency of task B
\item Give them the same name
\end{enumerate}
\end{frame}
\begin{frame}{Answer 96}
\textbf{Correct: C} --- dependencies define the DAG.\\
\vspace{0.5em}
\textbf{Wrong:}
\begin{itemize}
\item A: not a job feature.
\item B: not how it works.
\item D: not how it works.
\end{itemize}
\end{frame}

% ---------------------------------------------------------------
\begin{frame}{Question 97}
What does repairing a failed job run do?
\begin{enumerate}[label=\Alph*.]
\item Reruns only the failed and downstream tasks
\item Deletes the job
\item Reruns every task from scratch
\item Resets the cluster policy
\end{enumerate}
\end{frame}
\begin{frame}{Answer 97}
\textbf{Correct: A} --- successful tasks are not rerun.\\
\vspace{0.5em}
\textbf{Wrong:}
\begin{itemize}
\item B: not what it does.
\item C: that is a new run.
\item D: not what it does.
\end{itemize}
\end{frame}

% ---------------------------------------------------------------
\begin{frame}{Question 98}
How can a downstream task read a value produced by an upstream task?
\begin{enumerate}[label=\Alph*.]
\item Use \texttt{df.cache()}
\item Use task values with \texttt{dbutils.jobs.taskValues}
\item Use a global Python variable
\item Use \texttt{dbutils.fs.mv}
\end{enumerate}
\end{frame}
\begin{frame}{Answer 98}
\textbf{Correct: B} --- task values pass small values between tasks.\\
\vspace{0.5em}
\textbf{Wrong:}
\begin{itemize}
\item A: not shared across tasks.
\item C: tasks run in separate contexts.
\item D: moves files.
\end{itemize}
\end{frame}

% ---------------------------------------------------------------
\begin{frame}{Question 99}
Which statistical tests are commonly used to detect drift?
\begin{enumerate}[label=\Alph*.]
\item VIF and PCA
\item OPTIMIZE and VACUUM
\item ROC AUC and F1
\item Kolmogorov-Smirnov for numeric and chi-square for categorical features
\end{enumerate}
\end{frame}
\begin{frame}{Answer 99}
\textbf{Correct: D} --- they compare distributions.\\
\vspace{0.5em}
\textbf{Wrong:}
\begin{itemize}
\item A: not drift tests.
\item B: Delta maintenance.
\item C: model metrics.
\end{itemize}
\end{frame}

% ---------------------------------------------------------------
\begin{frame}{Question 100}
A model's online accuracy drops over time though code is unchanged. What is the best response?
\begin{enumerate}[label=\Alph*.]
\item Ignore it
\item Increase cluster size
\item Delete the model registry
\item Monitor for drift and retrain on fresh data when thresholds are exceeded
\end{enumerate}
\end{frame}
\begin{frame}{Answer 100}
\textbf{Correct: D} --- production data can change.\\
\vspace{0.5em}
\textbf{Wrong:}
\begin{itemize}
\item A: performance will continue to degrade.
\item B: does not fix drift.
\item C: does not help.
\end{itemize}
\end{frame}

\end{document}